% 二、研究方案

\section{研究目标、内容和拟解决的关键问题}

\subsection{研究目标}

建立一套基于基函数重构的软硬耦合问题求解理论框架，实现高效、鲁棒的物理仿真。具体目标包括：建立基函数重构的统一理论框架，阐明其与传统方法的本质区别；发展针对不同类型软硬耦合问题的专用基函数设计方法；实现近线性复杂度的高效求解算法；验证方法在典型应用场景中的有效性和实用性。

\subsection{研究内容}

本研究的核心思想是通过定制化的基函数与坐标变换，在新的表示子空间内"吸收"或"剥离"系统的硬约束与病态特征。

\textbf{研究内容一：不可拉伸弹性棒仿真（一维极度耦合）}

对于一维细长弹性结构（如毛发、绳索、缝合线等），其核心的软硬耦合体现在长度的绝对"不可伸长性"（无穷大拉伸刚度）与极小的弯曲/扭转刚度之间的矛盾。本研究提出紧凑坐标表示方法，通过一次深度的坐标基变换，摒弃传统的节点位置自由度，转而使用基于根节点位置以及局部轴角的曲率链式参数化体系。在这一全新的坐标基底下，系统的求解维度被精确缩减至与其内蕴自由度完全一致的$3n$维，所有硬约束在表示层面自动满足，将约束优化问题转化为无约束优化问题。

\textbf{研究内容二：大规模特征值问题求解（宏观全局耦合）}

在超大规模弹性系统中，系统的宏观低频柔性运动与各子结构界面处高频、强刚度的耦合约束交织在一起。本研究在CMS理论基础上，创新性地提出基于多重划分的相位互补模态基变换方法，将庞大的全节点物理空间映射至由不同空间划分策略提取出的降维模态子空间中，形成完美的"相位互补"。基于此，提出无Schur补及无局部特征求解的界面模态缩减法（SE-free IMR），直接利用对偶划分中的低频模态基近乎无损地逼近原始的界面约束空间。

\textbf{研究内容三：畸变单元稳定性处理（三维局部畸变）}

在三维弹性体有限元分析中，病态网格单元会产生极大的插值误差，导致该区域的人工刚度急剧增加。本研究提出谱几何退化分析，将系统的病态条件明确划分为"尺度病态"与"结构性病态"。为了消除结构性硬耦合，提出数值感知的局部细化与聚合策略，利用局部刚度矩阵的主特征向量作为物理探针，识别出产生最强人工锁定的方向，采用基于算子自适应的基函数优化与凝聚技术重新构造新插值基。

\subsection{拟解决的关键问题}

如何通过坐标变换消除不可拉伸弹性棒的内在约束？挑战在于需要在保持物理精度的同时将约束优化问题转化为无约束问题，解决思路是设计紧凑坐标表示使约束在表示层面自动满足。如何构造相位完备的模态基函数提升CMS收敛速度？挑战在于传统CMS子结构基函数相位固定，解决思路是利用多重划分生成相位互补的基函数集合。如何识别并过滤畸变单元产生的高频伪刚性模式？挑战在于传统网格细化方法代价高昂，解决思路是通过谱分析识别病态方向设计自适应基函数过滤。如何实现近线性复杂度的高效求解？解决思路是利用特殊矩阵结构设计高效预处理器。

\section{拟采取的研究方法、技术路线、试验方案及可行性分析}

\subsection{研究方法}

理论分析与数学推导方面，推导基函数变换的数学性质（正定性、完备性等），证明方法的收敛性和数值稳定性，分析计算复杂度和内存复杂度。数值实验与验证方面，设计标准测试算例，与现有方法进行定量对比，进行参数敏感性分析和鲁棒性测试。算法实现与系统集成方面，开发高效、模块化的算法实现，支持并行计算和GPU加速。

\subsection{技术路线}

\textbf{第一阶段：不可拉伸弹性棒仿真方法研究（2018.09 - 2020.08）}

任务包括：分析Cosserat杆的力学模型和约束条件，设计基于轴角表示的紧凑坐标参数化，推导紧凑表示下的运动学和动力学方程；分析系统矩阵的特殊结构（块三对角等），设计对称正定的混合预处理器；实现隐式时间积分器，设计主动集法处理碰撞约束，与Super-Helices、RedMax等方法进行对比实验。

\textbf{第二阶段：大规模特征值问题求解方法研究（2020.09 - 2021.08）}

任务包括：研究空间交错划分的构造方法，分析不同划分策略的相位互补特性，设计针对复杂拓扑的METIS划分集成方案；推导界面模态的近似公式，实现低频界面模态的高效计算；实现MPI并行算法，进行强扩展性和弱扩展性测试，在超算集群上验证大规模问题求解能力。

\textbf{第三阶段：畸变单元处理方法研究与论文撰写（2021.09 - 2022.06）}

任务包括：定义几何退化指标，建立尺度病态与结构性病态的分类准则，开发病态单元自动识别算法；利用局部刚度矩阵的主特征向量识别病态方向，设计自适应基函数优化策略；整合三种方法的共性理论，建立统一理论框架，完成博士学位论文撰写。

\subsection{试验方案}

试验平台包括桌面工作站（AMD Ryzen 9 7900X 12核处理器，128GB内存）和超算集群（Intel Xeon Gold 6248节点，每节点512GB内存）。测试用例包括标准弹性杆模型、二维/三维矩形区域模态分析、复杂几何模型（汽车部件、飞行器结构）和工业级大规模模型。评估指标包括计算时间、内存占用、特征值相对误差、迭代次数、收敛速度、扩展性效率等。对比方法包括传统有限元方法、Super-Helices方法、标准CMS方法、RedMax方法和商业软件（ANSYS、Abaqus）。

\subsection{可行性分析}

理论可行性方面，本研究基于成熟的数值分析和计算力学理论，基函数重构思想在谱方法、小波方法中有成功先例，初步实验结果已验证方法的有效性。技术可行性方面，研究团队具备丰富的算法开发经验，实验室拥有完善的软硬件设施，可利用开源库（如Eigen、PETSc等）加速开发。时间可行性方面，四年博士学制提供了充足的研究时间，三个研究内容相对独立可并行推进，已有初步研究成果。

\section{研究的创新点}

\subsection{理论创新}

提出"基函数重构"的新范式，区别于传统的代数预处理器方法，从离散表达层面解决软硬耦合问题。建立统一的理论框架，三个研究内容共享统一的理论内核：通过基函数变换将结构性病态转化为尺度病态或直接消除。揭示相位完备性的数学本质，在模态综合法研究中首次明确指出传统CMS收敛慢的根本原因是基函数相位固定。

\subsection{方法创新}

紧凑坐标表示方法将约束优化问题转化为无约束问题，系统维度从$7n-4$降至$3n$，所有内在约束在表示层面自然满足。近线性复杂度预处理器利用系统矩阵的特殊结构（块三对角），构造对称正定的混合预处理器，PCG收敛速度提升1-2个数量级。相位互补的多重划分策略利用Primal-Dual空间交错划分，生成相位完备的模态基函数，精度提升3个数量级。无Schur补的IMR方法避免了传统CMS中最昂贵的计算步骤，显著降低内存占用。算子自适应基函数基于局部刚度矩阵主特征向量识别病态方向，通过基函数重构过滤高频伪刚性模式。

\section{研究计划及预测进展}

\textbf{第一年（2018.09 - 2019.08）：文献调研与基础研究}

完成相关领域的文献调研，掌握有限元方法、模态综合法等基础理论，熟悉现有仿真软件和开发工具；深入研究不可拉伸弹性棒的力学模型，设计紧凑坐标表示方法，初步实现原型算法。预期成果：文献综述报告1份，初步算法原型。

\textbf{第二年（2019.09 - 2020.08）：不可拉伸弹性棒仿真方法研究}

完善紧凑坐标表示方法，设计混合预处理器，证明预处理器正定性；实现完整的隐式积分器，开发碰撞处理模块，进行大规模数值实验。预期成果：发表论文1篇。

\textbf{第三年（2020.09 - 2021.08）：大规模特征值问题求解方法研究}

设计多重划分策略，开发无Schur补的IMR方法，实现MPI并行算法；进行强扩展性和弱扩展性测试。预期成果：发表论文1篇。

\textbf{第四年（2021.09 - 2022.06）：理论整合与论文撰写}

整合三种方法的共性理论，建立统一理论框架，探索应用拓展方向；完成博士学位论文，准备答辩材料。预期成果：博士学位论文1篇，学术论文2篇。

\section{预期研究成果}

\subsection{学术成果}

预期发表高水平学术论文2篇，目标期刊/会议包括ACM Transactions on Graphics (TOG)、IEEE Transactions on Visualization and Computer Graphics (TVCG)、SIGGRAPH / SIGGRAPH Asia、Eurographics等。研究内容与论文规划：论文1为紧凑坐标表示方法与预处理器设计，论文2为多重划分模态综合法。

\subsection{人才培养}

完成博士学位论文1篇，培养独立科研能力；指导硕士生1-2名参与项目研究，培养团队协作能力。

\subsection{应用成果}

开发物理仿真原型系统，集成三种核心方法；与相关企业合作，探索技术转化可能性；在毛发动画、工业CAE软件中验证方法的有效性。